\section{Evidential relevance: complete derivations}
\label{sec:p15-r43-annex}

\begin{annexitem}{annex:p15.r43.polarity}{thm:p15-r43-polarity}{Polarity, conclusive chains, and the evidence application}
\begin{proof}
Symmetry follows from \(AB=BA\). Complement expansion gives
\(\Delta_P(\neg A,B)=-\Delta_P(A,B)\), and applying it twice gives
\(\Delta_P(\neg A,\neg B)=\Delta_P(A,B)\).
If \(P(A)=0\), then \(P(AB)=0\); if \(P(A)=1\), then
\(P(\neg A)=0\) and \(P(AB)=P(B)\). Both give zero covariance.
For \(P(A)>0\),
\(P(B\mid A)-P(B)=\Delta_P(A,B)/P(A)\).

If A is extremely positive to B and B extremely positive to C,
\(P(A\neg B)=P(B\neg C)=0\). Since
\(A\neg C\subseteq(A\neg B)\cup(B\neg C)\), its mass is zero.
Hence \(P(C\mid A)=1\), with \(P(A)>0,P(C)<1\), proving the chain.
Logical entailment gives \(A\subseteq B\), hence \(P(AB)=P(A)\)
and the claimed extreme positivity. Conversely, a probability model
can give a countercase \(A\neg B\) zero mass without removing that
possibility, so extreme relevance is not logical entailment.

In the three-outcome example the covariances of \((A,B),(B,C),(A,C)\)
are respectively \(1/9,1/9,-1/9\). This proves nontransitivity.
Finally, \(H\subseteq D\) gives \(P(H\mid D)=P(H)/P(D)\) and,
when \(P(D)<1\), \(P(H\mid\neg D)=0\). With \(P(D)=1\),
conditioning on D leaves every event probability unchanged.
\end{proof}
\end{annexitem}

\begin{annexitem}{annex:p15.r43.screening}{thm:p15-r43-screening}{The total-covariance calculation}
\begin{proof}
Let \(d_1=P(AC\mid B)-a_1c_1\) and
\(d_0=P(AC\mid\neg B)-a_0c_0\). Expand total probability for A,
C and AC and subtract their product to obtain
\[
 \Delta_P(A,C)=p d_1+(1-p)d_0+p(1-p)(a_1-a_0)(c_1-c_0).
\]
Also \(\Delta_P(A,B)=p(1-p)(a_1-a_0)\) and
\(\Delta_P(B,C)=p(1-p)(c_1-c_0)\). Conditional independence makes
\(d_0=d_1=0\), proving the identity and product rule for signs.
For positive links the product term is strictly positive, so
nonnegative \(d_0,d_1\) also suffice.

Under independence, \(P(AC\mid B)=a_1c_1\) and
\(P(AC\mid\neg B)=a_0c_0\). Positive links mean
\(a_1>a_0\geq0,c_1>c_0\geq0\), giving \(a_1c_1>a_0c_0\).
Thus \(\Delta_P(AC,B)>0\).

For the conclusive chain, C has conditional probability one within B,
so is independent of every event in that stratum. A has conditional
probability zero within \(\neg B\), because \(P(A\neg B)=0\), and
is likewise independent there. The strata are defined by \(0<p<1\).
The numerical group example has \(P(AC)=(16/25+1/25)/2=17/50\),
and division by \(P(A)=1/2\) gives \(17/25\).
\end{proof}
\end{annexitem}

\begin{annexitem}{annex:p15.r43.measures}{thm:p15-r43-measures}{The measures' identities and direction properties}
\begin{proof}
Write \(x=P(EH),e=P(E),h=P(H)\). On positive e and h,
\(q=x/e\), \(D=q-h\), and \(Q=x/(eh)\), proving the two
probability updates and Q symmetry. On \(0<h<1\),
\(a=x/h,b=(e-x)/(1-h)\), whence
\[
 L=\frac{x(1-h)}{h(e-x)}
   =\frac{q/(1-q)}{h/(1-h)}
\]
where denominators are positive. The zero/infinite cases follow from
one likelihood vanishing while the evidence has positive mass.
Taking logarithms where defined proves the additive identities.

For positive L,
\(\tanh(\frac12\ln L)=(L-1)/(L+1)=(a-b)/(a+b)\).
This transformation is strictly increasing, with limits -1 and 1
at zero and infinity. D and Q compare q to h; L compares a to b,
which has the same sign because
\(q-h=h(1-h)(a-b)/e\). KO has that numerator and a positive
denominator; Z divides the signed difference by a positive quantity.

Replacing H by its complement sends \(q-h\) to its negative, exchanges
a and b, and exchanges the positive/negative denominators of Z.
This proves target negation, including one-sided limits for L and LL.
Source negation gives
\(D(\neg E\Yright H)=-eD(E\Yright H)/(1-e)\), not generally -D.
For example, \(a=3/5,b=1/5\) gives \(L=3\), but the complementary source
has \(L=(2/5)/(4/5)=1/2\), not \(1/3\). Z likewise has no general source
inversion because its numerator is rescaled and its denominator changes
branch.

For A and B, putting \(a=P(A),b=P(B),x=P(AB)\), gives
the stated D magnitudes. When \(x>0\) and both \(x<a,b\),
\[
 L(A\Yright B)-L(B\Yright A)
 =\frac{x(b-a)(x-ab)}{ab(a-x)(b-x)}.
\]
For a positive covariance, the smaller source probability therefore
has the larger L; for negative covariance it has the smaller L,
further below one. If \(x=0\) both L values are zero. Positive extreme
cases follow by their infinite limits.

The ranges follow from \(0\leq q\leq1,0<h<1\) and the
displayed formulas: D cannot reach either signed unit, Q is finite
but unbounded, L can attain zero/infinity, and the logarithms have
their stated ranges. Nondegenerate h and positive evidence make
Z=1 exactly at q=1 and Z=-1 exactly at q=0; L has the same
conclusive endpoint condition. D increments of .2 are immediate.
The finite odds multiplier for .001 to .999 is \(999^2=998001\),
while an update from .999 to one has infinite posterior odds.
For Z the two affine branches meet at zero when \(q=h\). Their
slopes in q are \(1/h\) below h and \(1/(1-h)\) above h.
Both are positive, proving strict increase and continuity, but they
are equal only at \(h=1/2\), proving the differentiability assertion.
\end{proof}
\end{annexitem}

\begin{annexitem}{annex:p15.r43.sufficiency}{thm:p15-r43-sufficiency}{Finite binary minimal sufficiency}
\begin{proof}
Put \(s(x)=f_1(x)+f_0(x)>0\). The pair of ratios
\(f_1(x)/s(x),f_0(x)/s(x)\) is
\((\Lambda/(1+\Lambda),1/(1+\Lambda))\), with the endpoint
interpretation \((1,0)\) for infinite \(\Lambda\). Thus on a \(\Lambda\) fiber each
mass function is a constant multiple of the same \(s(x)\). Normalizing
within the fiber gives the same conditional distribution whenever the
corresponding fiber has positive mass. This proves sufficiency.

For any sufficient statistic \(T\), the common normalized data probabilities
within a fiber \(t\) give a factorization
\(f_i(x)=g_i(t)k_t(x)\), where \(g_i(t)\) is its fiber mass.
On supported points \(k_t(x)>0\), so the likelihood ratio is
\(g_1(t)/g_0(t)\), with the same infinite convention. It is constant
on the \(T\) fiber, proving that \(\Lambda\) is a function of \(T\). Logarithm is
injective on the extended nonnegative domain and so preserves fibers.
\end{proof}
\end{annexitem}

\begin{annexitem}{annex:p15.r43.characterization}{thm:p15-r43-characterization}{The complete realization argument for the ordinal families}
\begin{proof}
This is a construction proof following the dependencies' strategy,
with domain guards explicit. Write \(h=P(H),e=P(E),q=P(H\mid E)\).
The joint table is determined by \((h,e,q)\): its cells are
\(eq,e(1-q),h-eq,1-h-e(1-q)\), which must be nonnegative.
Formality therefore gives one function \(F(h,q,e)\) on the feasible triples.

\paragraph{Eliminate the evidence marginal.}
Fix \(h,q\) and two feasible \(e_i\). Within H define evidence events \(E_i\)
with probabilities \(e_iq/h\); within its complement use
\(e_i(1-q)/(1-h)\). These numbers lie in \([0,1]\) by feasibility.
Make \(E_1\),\(E_2\) conditionally independent in each stratum, using the
four products of their Bernoulli probabilities, weighted by \(h\) and \(1-h\).
This gives a finite model with both specified marginals and the same
posterior \(q\). POP-D equates their R values. Hence \(F(h,q,e)=G(h,q)\).
This step includes \(q=0,1\); zero
atom masses need not be interpreted as nonzero-support outcomes.

\paragraph{Difference family.}
Put \(d=q-h\). For two feasible h values \(h_1<h_2\) with the same d,
let \(g=h_2-h_1\). Construct disjoint targets \(H_1\) and J, with prior
masses \(h_1\) and \(g\). Given a small evidence event E, assign their
conditional masses \(h_1+d\) and \(g\), leaving \(1-h_2-d\) to the third
target cell. Outside E assign the masses obtained by subtracting
these evidence contributions from the priors and dividing by \(1-e\).
Choose \(e>0\) sufficiently small that every cell is nonnegative:
\(e(h_1+d)\leq h_1\) and
\(e(1-h_2-d)\leq1-h_2\); the J inequality is automatic.
The zero numerator cases need no bound. Because \(h_2<1\) and \(h_1>0\),
such an \(e<1\) exists. J has unchanged posterior \(g\), and MDHD therefore
equates the R values for \(H_1\) and \(H_1\) union J. Consequently
\(G(h,h+d)=f_D(d)\), independent of \(h\).

For \(d_1>d_2\) differing by less than one, choose
\(h\in(0,1)\) with \(0\leq h+d_2<h+d_1\leq1\).
The evidence construction above realizes both posteriors in one model;
POP-D makes \(f_D(d_1)>f_D(d_2)\). If the difference is at least one,
then \(d_1>0>d_2\), and comparisons with zero give the same strict
order. Thus \(f_D\) is strictly increasing on \((-1,1)\). Conversely,
such a transform obeys POP-D and MDHD since the disjoint-target
difference satisfies \(D(E,H_1\cup H_2)-D(E,H_1)=D(E,H_2)\).

\paragraph{Quotient family.}
Fix \(r=q/h\) and two feasible \(h_i\), so \(q_i=rh_i\) lies in \([0,1]\).
Choose \(e\in(0,1)\) small enough that
\(e q_i\leq h_i\) and \(e(1-q_i)\leq1-h_i\) for both i.
Within E, construct \(H_1\),\(H_2\) independently with marginal probabilities
\(q_i\); outside E use marginal probabilities
\((h_i-eq_i)/(1-e)\) and again their Bernoulli product.
This realizes both target tables in one finite model.
Their likelihoods \(P(E\mid H_i)=er\) are equal, so LIK-D equates
their R values. Hence \(G(h,q)=f_Q(q/h)\).
For any \(r_1<r_2\) choose one positive \(h\) small enough that both
\(r_i h\) are at most one, then choose small evidence masses and use
the first construction. POP-D makes \(f_Q\) strictly increasing.
Conversely \(q/h\) is increasing in \(q\) for fixed \(h\), and equals
\(P(E\mid H)/P(E)\) for a fixed source E, proving both axioms.

\paragraph{Likelihood family.}
Fix a likelihood ratio \(\lambda\) and two priors \(h,h'\in(0,1)\).
Choose \(a,b\in[0,1]\) with extended ratio \(a/b=\lambda\), not both zero, and small
enough that the evidence event \(E_2\) is not certain. For finite positive
\(\lambda\) one may take \(b>0\) small and \(a=\lambda b\); for zero use \(a=0\),
and for infinity use \(b=0,a>0\). Choose an \(E_1\) with positive
likelihoods \(u,v\) whose ratio is \(O(h')/O(h)\), scaling both sufficiently
small to keep them below one. Its posterior target probability is \(h'\).
Make \(E_1\),\(E_2\) independent given H and given its complement by
Bernoulli products. MCIE equates the relevance of \(E_2\) before and
after \(E_1\), while its likelihoods \(a,b\) stay unchanged. Thus \(G(h,q)\)
depends only on \(\lambda=O(q)/O(h)\). With \(h\) fixed and finite
\(\lambda\), the posterior is \(q=\lambda h/(\lambda h+1-h)\).
At \(\lambda=+\infty\), the case \(b=0<a\) gives \(q=1\),
rather than evaluating an infinite-over-infinite fraction. These values
increase strictly from zero to one on the extended ratio domain.
POP-D proves strict increase of \(f_L\), including its ordered endpoints.
Conversely the odds formula proves
POP-D, and conditional independence preserves a and b, proving MCIE.

All constructions use finitely many worlds with nonnegative masses
summing to one. Comparisons that involve a certain/impossible target
or an undefined conditioning event were explicitly excluded from the
axiomatic domain; no neutrality completion is used to bridge them.
The argument needs no smoothness or numerical continuum of worlds.
\end{proof}
\end{annexitem}

\begin{annexitem}{annex:p15.r43.rankings}{thm:p15-r43-rankings}{Ranking proofs, realizations, and the raven calculation}
\begin{proof}
POP-D supplies exactly the same comparisons with the same posterior
schedule for every measure satisfying it. For D,Q,L this also follows
from q-h,q/h and O(q)/O(h), each increasing in q with fixed h.
Z is increasing on each side of q=h, meets continuously at zero,
and has positive slope on both sides. Increasing transforms preserve
these comparisons and ties.

For the two target tables, substitute the displayed h,q into those
three expressions to obtain their D,Q,L columns. Realizability is not
assumed: given E of mass 1/5, make \(H_1\),\(H_2\) conditionally independent
inside and outside E, with inside marginals \(q_i\) and outside marginals
\((h_i-q_i/5)/(4/5)\). In model I these outside marginals are
21/40 and 1/20; in model II they are 1/8 and 61/80. All lie in
[0,1], so Bernoulli products define the required eight-world laws.
The inequalities in the table prove the three pairwise divergences.

For the direction counterexample, within H give \(E_1\),\(E_2\) independent
probabilities 2/25 and 7/10, and within its complement give them
1/50 and 3/10. Weight the two strata by 1/2. These give evidence
marginals 1/20,1/2 and posteriors 4/5,7/10 as stated.
Reverse D is P(\(E_i\)|H)-P(\(E_i\)), hence 3/100 and 1/5.
Q symmetry was already proved. Thus Q can represent the robust
forward ordering nondirectionally, while arbitrary reversal of D cannot.

For the raven model, conditional on H the four (raven,nonblack)
cells have masses \(0,1/100,4/5,19/100\); conditional on its
complement they have \(1/200,1/200,159/200,39/200\), in the same
order: both, raven only, nonblack only, neither. Each table sums to
one and has the stipulated marginals. The black-raven event has
likelihoods 1/100 and 1/200, so its equal-prior posterior is 2/3
and L is 2. The nonblack-nonraven event has likelihoods 4/5 and
159/200, giving posterior 160/319 and L=160/159. Both exceed
the prior 1/2, with the claimed ordering.
If a shoe event C is a subset of that B event and is conditionally
independent of H within B, with positive probability, then
\(P(H\mid C)=P(H\mid B)\), by the conditional independence formula. \(C\subseteq B\)
alone imposes no such equality. A single nonblack nonraven specimen
also does not logically entail the universal contrapositive of H.

For the Linda-structured numerical model in the body, multiplying
the E-conditional cells by \(P(E)=1/10\) gives
\((.01,.005,.07,.015)\). Subtraction from the prior cells gives
the joint masses outside E, \((.01,.175,.41,.305)\), all nonnegative
and summing to .9. This constructs one normalized eight-world law.
The conjunction has prior .02 and posterior .10, giving D=.08;
bank-teller work has prior .20 and posterior .15, giving D=-.05;
feminist activity has prior .50 and posterior .80. Inclusion is
unchanged, since \(.10\leq.15\).
\end{proof}
\end{annexitem}

\begin{annexitem}{annex:p15.r43.composition}{thm:p15-r43-composition}{Conjunctive products, disjunctive mixtures, and exclusive or}
\begin{proof}
Conditional independence factors each conjunction likelihood into its
two factors; dividing gives the L product. Logarithm gives the LL sum
when all factors are positive and finite. KO, for example, need not
turn that product into either a sum or a product of KO values.

For disjoint sources, total conditional probability gives
\[
 P(H|A\cup B)=\frac{P(A)P(H|A)+P(B)P(H|B)}{P(A)+P(B)}.
\]
Subtracting h or dividing by h gives the D/Q means. Null disjuncts
have weight zero and can be omitted, so no ratio on a null disjunct
is needed. For L use
\(P(A\cup B|H)=P(A|H)+P(B|H)\) and the same identity under
the complement. A ratio of the sums is the weighted mean of ratios,
weighted by the denominator likelihoods. The hypothesis and source
support conditions ensure that all these denominators are defined.
Increasing transforms put the transformed mean between transformed
endpoints; they do not commute with the original weighted average.

Indicators satisfy \(1_{A\cup B}=1_A+1_B\) almost surely for
disjoint A,B, and \(1_{AB}=1_A+1_B-1\) almost surely if their union
has mass one. Taking covariance with \(1_H\) proves the Carnap
claims. Since \(R_C(E,H)=P(E)D(E,H)\), it fails POP-D: the
direction-counterexample sources have forward D values .3,.2, but
covariances .015,.1, with the opposite ranking.

For the final configuration let \(a_i=P(E_i\mid H)\) and
\(b_i=P(E_i\mid\neg H)\). Since \(0<b_i<a_i\leq1\),
the finite ratios \(l_i=a_i/b_i\) exceed one.
Conditional independence gives conjunction ratio \(l_1l_2\),
greater than either. Union likelihoods are
\(S-A\) and \(T-B\), where
\(S=a_1+a_2,T=b_1+b_2,A=a_1a_2,B=b_1b_2\).
Their difference is
\((a_1-b_1)(1-a_2)+(a_2-b_2)(1-b_1)>0\), so \(L_{\cup}>1\).
Let \(l_1\) be the larger ratio. Then
\[
 l_1(T-B)-(S-A)
 =(l_1-l_2)b_2+l_1(l_2-1)b_1b_2>0,
\]
giving \(L_{\cup}<\max(l_1,l_2)\).

Exclusive-or likelihoods are \(S-2A\) and \(T-2B\). The latter is positive
because \(0<b_i<1\). Comparing its ratio with the union ratio has
cross-product numerator \(SB-AT\), which is strictly negative:
\(S/A=1/a_1+1/a_2<1/b_1+1/b_2=T/B\).
For symmetric likelihood pairs \((a_i,b_i)=(2/5,1/5)\),
\((3/4,1/4)\), and \((4/5,2/5)\), the exclusive-or ratios are respectively
\(3/2,1,2/3\). Thus every polarity is realized under the premises.
Finally, total probability gives
\(P(H)=eP(H\mid E)+(1-e)P(H\mid\neg E)\), whence
\(P(H\mid E)-P(H\mid\neg E)=D(E,H)/(1-e)\) for \(0<e<1\).
\end{proof}
\end{annexitem}
